\documentclass[10pt,a4paper]{amsart}
\usepackage[margin=19mm]{geometry}
\usepackage{amsmath,amssymb,mathtools}
\usepackage[scr=rsfs,cal=euler]{mathalfa}
\usepackage{xfrac}
\usepackage[unicode,hidelinks,bookmarks=false]{hyperref}
\newcommand{\ie}{i.e.\ }
\newcommand{\cf}{cf.\ }
\newcommand{\resp}{respectively\ }
\newcommand{\Subsection}{\subsection}
\providecommand{\nobreakdash}{\nobreak\hbox{-}}
\setlength{\emergencystretch}{2em}
\multlinegap0pt
\usepackage[shortlabels]{enumitem}
\usepackage{amssymb}
\usepackage{mathtools}
\usepackage{bm}
\usepackage{tikz}
\newtheorem{thm}{Theorem}
\title{Splitting fields and spectral invariants of character degree graphs in solvable groups}
\author{G. Sivanesan, C. Selvaraj, Jacob Laubacher}
\date{Source: arXiv:2608.14304v1 (CC BY 4.0)}
\begin{document}
\maketitle

\noindent\emph{Source and licence.} Splitting fields and spectral invariants of character degree graphs in solvable groups. Version arXiv:2608.14304v1 (CC BY 4.0), distributed by arXiv under the Creative Commons Attribution 4.0 International licence, http://creativecommons.org/licenses/by/4.0/, as recorded in the arXiv licence metadata for this exact version and shown on the abstract page https://arxiv.org/abs/2608.14304v1. Full licence notice: \texttt{licenses/arxiv-2608.14304-CC-BY-4.0.txt}.\par\smallskip

\noindent\textit{Editorial scope.} Counted unit: the article's thm thm:LewisSpectralCorrected (source0.tex lines 744--773). The article's complete original proof of it is delivered with it (source0.tex lines 775--867, one \texttt{proof} environment of 93 lines). No mathematics is written, repaired or re-derived: the statement and the proof are the article's own lines, in the article's own notation, and no retained line is substituted or re-flowed.  No further source-local context is needed: every cross-reference of every retained line resolves inside the two delivered spans.  Nothing was omitted from any retained span, so the package carries no omission record.  No retained line contains a citation, so the wrapper carries no bibliography and no bibliography entry.  No retained line contains a figure environment, an image-inclusion command or drawing code, so no image is delivered and the package declares no asset dependency.  Editorial formatting only, in addition: an AMS \texttt{amsart} preamble that restates the article's own declarations and macros used by the retained lines, namely the environments \texttt{thm} on separate counters, together with the standard packages those lines need.  The retained environments print in this document as Theorem 1 in order of appearance, whereas the article prints the same items with the numbering of its own full document.  The complete native source of the article as distributed in the arXiv e-print for this exact version is bundled unchanged as \texttt{sources/207656/source0.tex}.
\begin{thm}[Spectral Structure of Lewis Graphs]\label{thm:LewisSpectralCorrected}
Let $\Gamma$ be a Lewis graph with vertex partition
\[
V(\Gamma)=U\sqcup X\sqcup Y, \qquad |U|=n_1,\; |X|=n_2,\; |Y|=n_3,
\]
satisfying the following structural conditions:
\begin{enumerate}[label=\upshape(\roman*)]
    \item The induced subgraphs $\Gamma[U]$, $\Gamma[X]$, and $\Gamma[Y]$ are complete;
    \item Every vertex $u\in U$ is adjacent to all vertices in $X\cup Y$;
    \item For each $x\in X$, there is exactly one $y\in Y$ such that $xy\notin E(\Gamma)$, and this assignment is given by a surjective map $f\colon X\to Y$;
    \item Let $B_{XY}\in\{0,1\}^{n_2\times n_3}$ denote the \emph{bipartite defect matrix}, defined by $(B_{XY})_{x,y}=1$ if $xy\notin E(\Gamma)$ and $0$ otherwise, and let $r=\operatorname{rank}(B_{XY})$ over $\mathbb{Q}$.
\end{enumerate}
Let $A=A(\Gamma)$ be the adjacency matrix of $\Gamma$, and let $n = n_1+n_2+n_3$. Then the following hold:
\begin{enumerate}[label=\upshape(\alph*)]
    \item \textbf{(Guaranteed integer eigenvalues).}
    The eigenvalue $-1$ occurs with multiplicity at least
    \[
      \max\{0,\, n - 1 - 2r\}.
    \]
    \item \textbf{(Bound on exceptional eigenvalues).}
    At most $2r+1$ eigenvalues of $A$ differ from $-1$ (counting algebraic multiplicities).
    \item \textbf{(Degree of irreducible factors).}
    Every irreducible factor of the characteristic polynomial $\chi_A(\lambda)\in\mathbb{Q}[\lambda]$ corresponding to eigenvalues not equal to $-1$ has degree at most $2r+1$.
    \item \textbf{(Splitting-field bound).}
    Let $K$ be the splitting field of $\chi_A(\lambda)$ over $\mathbb{Q}$. Then
    \[
      [K:\mathbb{Q}] \leq (2r+1)^{\,2r+1}.
    \]
\end{enumerate}
\end{thm}

\begin{proof}
We proceed in four steps.

\medskip
\noindent\textbf{Step 1: Decomposition $A = A_0 + E$.}

Let $A_0$ be the adjacency matrix of the complete graph $K_n$, obtained from $\Gamma$ by adding all missing edges between $X$ and $Y$. Then
\[
A_0 = J_n - I_n,
\]
where $J_n$ is the all-ones matrix and $I_n$ is the identity. The spectrum of $A_0$ is
\[
\sigma(A_0) = \{\,n-1,\, -1^{(n-1)}\,\},
\]
so $A_0$ has eigenvalue $n-1$ with multiplicity $1$ and eigenvalue $-1$ with multiplicity $n-1$.

Define the defect matrix $E = A - A_0$. Ordering vertices as $(U,X,Y)$, we have
\[
E =
\begin{pmatrix}
0 & 0 & 0 \\
0 & 0 & -B_{XY} \\
0 & -B_{XY}^\top & 0
\end{pmatrix}.
\]
Since $B_{XY}$ has rank $r$, the block matrix $\begin{pmatrix} 0 & B_{XY} \\ B_{XY}^\top & 0 \end{pmatrix}$ has rank $2r$. Therefore
\[
\operatorname{rank}(E) = 2r.
\]

\medskip
\noindent\textbf{Step 2: The $-1$ eigenspace (Proof of part~\textup{(a)}).}

Let $\mathbf{1}\in\mathbb{R}^n$ denote the all-ones vector. Since $A_0 = J_n - I_n$, we have
\[
A_0 \mathbf{1} = (n-1)\mathbf{1}, \qquad \text{and} \qquad A_0 v = -v \;\text{ for all } v \perp \mathbf{1}.
\]
Define the subspace
\[
W = \mathbf{1}^\perp \cap \ker(E) \subseteq \mathbb{R}^n.
\]
For any $v \in W$, we have $v \perp \mathbf{1}$ and $Ev = 0$, so
\[
Av = (A_0 + E)v = A_0 v + 0 = -v.
\]
Thus every vector in $W$ is an eigenvector of $A$ with eigenvalue $-1$.

By the rank-nullity theorem, $\dim(\ker(E)) = n - \operatorname{rank}(E) = n - 2r$. Since $\mathbf{1}^\perp$ has codimension $1$ in $\mathbb{R}^n$, we obtain
\[
\dim(W) \geq \dim(\ker(E)) - 1 = (n - 2r) - 1 = n - 1 - 2r.
\]
Since multiplicities are non-negative, the eigenvalue $-1$ occurs with multiplicity at least $\max\{0, n - 1 - 2r\}$, proving part~\textup{(a)}.

\medskip
\noindent\textbf{Step 3: The exceptional subspace (Proof of parts~\textup{(b)} and~\textup{(c)}).}

Let $S = \operatorname{span}\{\mathbf{1}\} + \operatorname{im}(E)$. Since $\operatorname{im}(E)$ has dimension $\operatorname{rank}(E) = 2r$, we have
\[
\dim(S) \leq 1 + 2r.
\]
We claim that $S$ is $A$-invariant. Indeed:
\begin{itemize}
    \item $A\mathbf{1} = A_0\mathbf{1} + E\mathbf{1} = (n-1)\mathbf{1} + E\mathbf{1} \in S$, since $E\mathbf{1} \in \operatorname{im}(E) \subseteq S$;
    \item For any $w \in \operatorname{im}(E)$, write $w = Ez$. Then
    \[
    Aw = (A_0 + E)Ez = A_0 Ez + E^2 z.
    \]
    Since $A_0 = J_n - I_n$ and $J_n E = \mathbf{1}(\mathbf{1}^\top E)$, we have $A_0 Ez \in \operatorname{span}\{\mathbf{1}\} + \operatorname{im}(E) = S$. Also $E^2 z \in \operatorname{im}(E) \subseteq S$. Thus $Aw \in S$.
\end{itemize}
Therefore $S$ is $A$-invariant, and $\mathbb{R}^n = S \oplus S^\perp$ is an $A$-invariant orthogonal decomposition.

On $S^\perp$, we have $S^\perp \subseteq \mathbf{1}^\perp \cap \ker(E) = W$, so $A|_{S^\perp} = -I$. Hence all eigenvalues of $A$ not equal to $-1$ arise from the restriction $A|_S$, which is a matrix of size at most $(2r+1) \times (2r+1)$.

This proves:
\begin{itemize}
    \item At most $\dim(S) \leq 2r+1$ eigenvalues differ from $-1$ (part~\textup{(b)});
    \item Every irreducible factor of $\chi_A(\lambda)$ corresponding to eigenvalues $\neq -1$ divides the characteristic polynomial of $A|_S$, hence has degree at most $2r+1$ (part~\textup{(c)}).
\end{itemize}

\medskip
\noindent\textbf{Step 4: Splitting field bound (Proof of part~\textup{(d)}).}

Let $p(\lambda) \in \mathbb{Q}[\lambda]$ be the characteristic polynomial of $A|_S$. Then $\deg(p) = d \leq 2r+1$, and the splitting field $K$ of $\chi_A(\lambda)$ coincides with the splitting field of $p(\lambda)$ (since all other roots are $-1$).

The Galois group of $p(\lambda)$ over $\mathbb{Q}$ embeds into the symmetric group $S_d$, so
\[
[K:\mathbb{Q}] \leq |S_d| = d! \leq d^{\,d} \leq (2r+1)^{\,2r+1},
\]
proving part~\textup{(d)}.

\medskip
This completes the proof of Theorem~\ref{thm:LewisSpectralCorrected}. \hfill $\square$
\end{proof}

\end{document}
